\chapter{Brain Tumor Treatment}

\section*{Definition and Overview}
Treatment of a brain tumour refers to the coordinated medical and surgical approaches used to control or cure tumours growing in or around the brain. The choice of treatment depends on whether the tumour is benign or malignant, its specific type and grade, its size and location, and the person's age and overall fitness. The aims of treatment are graduated: to remove the tumour completely where possible; to slow its growth, relieve symptoms, and preserve neurological function where it cannot be removed; and, where cure is not possible, to maintain the best possible quality of life for as long as possible. Most people with a brain tumour receive a combination of treatments -- typically surgery followed by radiotherapy and/or chemotherapy -- planned and delivered by a multidisciplinary team. Treatment is increasingly individualised according to the molecular characteristics of each tumour.

\section*{Epidemiology}
Brain tumour treatment is a substantial and growing field of medicine. Primary brain tumours are comparatively uncommon, but metastatic brain tumours are frequent and increasing as people live longer with systemic cancers, so large numbers of patients are treated each year. The pattern of treatment mirrors the tumour mix: many benign tumours, such as meningiomas, are treated with surgery alone and never require further therapy, whereas malignant gliomas almost always need combined surgery, radiotherapy, and chemotherapy. Advances in imaging, surgical technique, radiotherapy delivery, and systemic therapy have steadily improved survival and quality of life, with the largest gains seen in younger, fitter patients and in tumour types with favourable molecular profiles.

\section*{Aetiology and Pathophysiology}
Treatments are selected according to how a tumour behaves and grows:

\begin{itemize}
    \item \textbf{Benign tumours:} such as meningiomas, grow slowly, do not spread beyond the nervous system, and are frequently cured by complete surgical removal
    \item \textbf{Malignant tumours:} such as glioblastoma, grow rapidly, invade surrounding brain tissue, and recur despite treatment
    \item \textbf{Metastatic tumours:} have spread from a cancer elsewhere and are treated in the context of the underlying disease, often in addition to systemic therapy for the primary cancer
    \item \textbf{Molecular characteristics:} genetic markers such as IDH mutation status and MGMT promoter methylation help predict behaviour and response to chemotherapy and targeted agents, and increasingly define treatment
\end{itemize}

The mechanisms of treatment mirror these differences: surgery physically removes tumour tissue; radiotherapy and chemotherapy damage the DNA of dividing tumour cells; targeted drugs block specific growth signals; and corticosteroids rapidly reduce the oedema and pressure around a tumour, providing symptom relief while definitive treatment is arranged.

\section*{Clinical Features}
Before treatment, symptoms reflect the tumour itself -- headache, seizures, weakness, speech difficulty, or personality change. Treatment itself also produces symptoms, which vary with the modality:

\begin{itemize}
    \item \textbf{After surgery:} headache, wound discomfort, drowsiness, and swelling of the scalp and face, usually settling over days to weeks
    \item \textbf{During radiotherapy:} fatigue, scalp irritation, hair loss over the treated area, and sometimes a temporary worsening of neurological symptoms from radiation-related brain swelling
    \item \textbf{During chemotherapy:} fatigue, nausea, reduced appetite, and an increased risk of infection from lowered blood counts
    \item \textbf{With corticosteroids:} increased appetite, weight gain, fluid retention, mood change, and difficulty sleeping
\end{itemize}

\section*{Differential Diagnosis}
New symptoms during treatment must be assigned to their correct cause:

\begin{itemize}
    \item Tumour progression or recurrence versus treatment-related side effects
    \item Radiation necrosis (damage to healthy brain after radiotherapy) versus recurrent tumour, which can look very similar on imaging
    \item Infection, including wound infection or meningitis after surgery, versus sterile inflammatory change
    \item Seizures from the tumour or after surgery versus other causes of collapse
    \item Corticosteroid complications such as hyperglycaemia or mood disturbance versus worsening of the underlying condition
    \item Raised intracranial pressure from fluid accumulation or hydrocephalus versus simple postoperative headache
\end{itemize}

\section*{When to Seek Medical Care}
People receiving brain tumour treatment should contact their team for persistent fever, worsening headache, new weakness or confusion, seizures, or any problem with the surgical wound. Many side effects are manageable with simple measures, but some require urgent review, and early treatment of complications protects quality of life and outcome.

\textbf{Call 000 (or local emergency services) for:}
\begin{itemize}
    \item A seizure or fitting episode
    \item A sudden, severe headache or sudden worsening of symptoms
    \item Confusion, drowsiness, or loss of consciousness
    \item New weakness, numbness, or difficulty speaking
    \item Fever with rigors or feeling generally very unwell, especially after recent chemotherapy that has lowered immunity
\end{itemize}

\section*{Diagnosis and Investigations}
\begin{itemize}
    \item \textbf{Imaging:} MRI and CT before, during, and after treatment to plan surgery and radiotherapy, monitor response, and detect complications such as hydrocephalus
    \item \textbf{Biopsy and pathology:} histological analysis of tumour tissue to establish the exact diagnosis and grade, which underpins every treatment decision
    \item \textbf{Molecular testing:} analysis of markers such as IDH status, MGMT promoter methylation, and 1p/19q codeletion, which guide the choice and intensity of treatment
    \item \textbf{Blood tests:} checks of blood counts and organ function before and during chemotherapy, and of glucose during corticosteroid use
    \item \textbf{Assessment of fitness:} evaluation of overall health, daily functioning, and neurological status, which helps calibrate how aggressive treatment should be
    \item \textbf{Monitoring scans:} scheduled imaging during follow-up to detect recurrence as early as possible
\end{itemize}

\section*{Management}
Treatment is individualised by a multidisciplinary team and typically combines several approaches:

\begin{itemize}
    \item \textbf{Surgery:} maximal safe resection of the tumour, or a diagnostic biopsy where resection would endanger vital brain areas; surgery may also relieve raised intracranial pressure quickly
    \item \textbf{Radiotherapy:} external beam radiation to the tumour and a surrounding margin, often delivered in daily fractions over several weeks; stereotactic radiosurgery gives a high, focused dose to small or inaccessible tumours in one or a few sessions
    \item \textbf{Chemotherapy:} drugs such as temozolomide for gliomas, given concurrently with and after radiotherapy; other regimens are used for lymphomas, metastases, and paediatric tumours
    \item \textbf{Targeted and immunotherapy:} newer agents directed at specific molecular features of the tumour, used increasingly for selected tumour types
    \item \textbf{Supportive medicines:} corticosteroids to reduce brain swelling, antiepileptic drugs to control seizures, and antiemetics and analgesia for symptom relief
    \item \textbf{Rehabilitation:} physiotherapy, occupational therapy, speech and language therapy, and neuropsychological support during and after treatment
    \item \textbf{Palliative care:} expert symptom control, psychological support, and advance care planning for tumours that cannot be cured
    \item \textbf{Clinical trials:} access to investigational treatments, discussed with the team where appropriate
\end{itemize}

\section*{Complications}
\begin{itemize}
    \item Surgical complications: bleeding, infection, brain swelling, and new neurological deficits
    \item Radiation side effects: early fatigue and hair loss, and later radiation necrosis of healthy brain tissue
    \item Chemotherapy side effects: fatigue, nausea, low blood counts, increased infection risk, and occasional effects on other organs
    \item Corticosteroid complications: weight gain, high blood sugar, osteoporosis, and mood disturbance with prolonged use
    \item Seizures despite antiepileptic medication
    \item Tumour recurrence after apparently successful treatment
    \item Psychological distress and the substantial practical and emotional burden of treatment on the person and their family
\end{itemize}

\section*{Prognosis and Outcomes}
Outcomes depend heavily on tumour type and molecular profile. Some benign tumours are cured by surgery alone, and many people with low-grade tumours live for years with excellent quality of life. Malignant gliomas such as glioblastoma are generally not curable, but combined treatment has meaningfully extended survival and a small minority of patients do considerably better than average. Metastatic brain tumours respond well in some people, particularly when the underlying cancer is responsive to systemic therapy. Because both the tumour and its treatment can affect thinking, energy, and mood, rehabilitation, psychological support, and careful symptom management are central to the best possible outcome, alongside the treatment itself.

\section*{Prevention and Self-Care}
\begin{itemize}
    \item Follow the treatment plan agreed with your team, and raise concerns rather than stopping or altering treatment on your own
    \item Take prescribed medicines -- corticosteroids, antiepileptic drugs, and antiemetics -- exactly as directed
    \item Maintain good nutrition and hydration, and balance rest against gentle activity to manage fatigue
    \item Protect the scalp and skin during radiotherapy as advised, and keep the surgical wound clean and dry
    \item Attend all follow-up scans and appointments, and report new or worsening symptoms promptly
    \item Arrange practical and emotional support, since a brain tumour affects work, family, and daily life
    \item If the tumour cannot be cured, discuss advance care planning and your preferences with the team so that future care reflects your wishes
\end{itemize}

\section*{Sources and Further Reading}
The following organisations provide reliable, evidence-based information for patients, students, and clinicians.

\begin{itemize}
    \item NHS. \textit{Brain tumour treatment: surgery, radiotherapy, and chemotherapy information.} https://www.nhs.uk/conditions/
    \item Mayo Clinic. \textit{Brain tumor treatment options and patient education.} https://www.mayoclinic.org/diseases-conditions/
    \item National Institute for Health and Care Excellence. \textit{Clinical guidance on brain tumour treatment pathways.} https://www.nice.org.uk/
    \item MSD Manual Professional Version. \textit{Brain tumor treatment: principles and modalities.} https://www.msdmanuals.com/
    \item MedlinePlus. \textit{Brain tumor treatments and supportive care: consumer health information.} https://medlineplus.gov/
    \item Centers for Disease Control and Prevention. \textit{Cancer treatment and survivorship information.} https://www.cdc.gov/
\end{itemize}
